% ============================================================
%  APPENDIX TO CHAPTER 4.  Material NOT in Rudin, imported from
%  Tao I/II, ECNU, USTC, Yu Pin.  Numbered S4.n throughout.
% ============================================================
\rappendix{4}

\begin{roadmap}[How to use this appendix]
Nothing in this appendix is Rudin's text. It is imported from the four
comparison texts and set apart at every level: the running head names the
appendix, section heads carry a reverse-video tag, item numbers are prefixed
\textbf{S}, and every statement names its source.

\medskip
\raggedright
\textbf{Why import anything at all.} Rudin's Chapter 4 is a chapter about
\emph{structure}: what continuity does to compactness and connectedness. He
buys that generality by postponing every concrete function --- $\sin$, $\exp$,
$\log$, $x^\alpha$ do not exist until Chapter 8 --- and by leaving the
working criteria for limits and uniform continuity to the exercises. The
Chinese and Tao texts make the opposite trade. Reading them alongside
supplies five things Rudin's chapter genuinely lacks:

\medskip
\textbf{1. Criteria (A).} A Cauchy criterion for function limits; one
definition of limit that specialises to all of Rudin's; a squeeze law; and
oscillation, which measures \emph{how} discontinuous a function is at a point.

\medskip
\textbf{2. Asymptotics (B).} The $o$/$O$ calculus and equivalent
infinitesimals --- the machinery every later computation runs on, and which
Rudin never introduces at all.

\medskip
\textbf{3. The postponed limits (C).} $\lim_{x\to0}(\sin x)/x = 1$ and
$\lim_{x\to\infty}(1+1/x)^x = e$, with the proofs the other books give.

\medskip
\textbf{4. Tools left in Rudin's exercises (D, E).} Inverse functions without
compactness; the sequential criterion for uniform continuity; gluing;
extension to the closure. Rudin asks for several of these as Exercises
4.11--4.13; the other texts prove them in the running text and then use them.

\medskip
\textbf{5. Finer structure (F).} Rudin's two kinds of discontinuity refine to
four; connectedness gains a path-connected cousin and a decomposition into
components.

\medskip
\textbf{How to read it.} None of this is needed to follow Rudin, which is why
it sits in an appendix rather than in his chapter. Sections C and D in
particular are \emph{logically} out of order in his development --- they use
functions he has not built yet --- and each says so where it matters. Read
the appendix as a second pass over Chapter 4, or as the answer to ``how would
I actually compute with this?'' The exercises begin on the last section.
\end{roadmap}

% ============================================================
\ssection{Supplement A --- Criteria for limits that Rudin omits}

\begin{roadmap}[Why this supplement belongs to Chapter 4]
\raggedright
\textbf{What Rudin does.} Definition 4.1 defines $\lim_{x\to p}f(x)=q$ by
$\varepsilon$-$\delta$, and 4.2 characterises it by sequences: $f(p_n)\to q$
for every sequence $p_n \to p$ in $E$ with $p_n \ne p$.

\smallskip
\textbf{The gap.} Both forms presuppose the candidate $q$. Neither can be
used to decide whether a limit \emph{exists} when no value suggests itself,
and Rudin supplies no criterion that does --- no Cauchy condition, no
squeeze, and no way of measuring how badly a limit fails to exist.

\smallskip
\textbf{Where the gap bites.} At Exercise 4.13, where a uniformly continuous
function must be extended to the closure and the limit at each new point has
to be produced from nothing; Rudin's hint amounts to proving a Cauchy
criterion by hand in that special case. It bites again at 4.26--4.28, where
discontinuities are classified but never quantified, so ``how discontinuous''
has no answer.

\smallskip
\textbf{What this section supplies.} The Cauchy criterion for function
limits; one definition of limit that specialises to all four of Rudin's; the
squeeze law he uses but never states; and oscillation, which turns
``continuous at $p$'' into the numerical statement $\omega_f(p)=0$ and makes
the set of continuity points a computable object.
\end{roadmap}

\sitem{S4.1}{Theorem}
\emph{(Cauchy criterion for function limits.) Let $X$ be a metric space, $E
\subset X$, $p$ a limit point of $E$, and $f : E \to Y$ with $Y$ complete
(for instance $Y = \R^k$). Then $\lim_{x\to p} f(x)$ exists if and only if:
for every $\varepsilon > 0$ there is a $\delta > 0$ such that}
\[ d_Y\bigl(f(x'),f(x'')\bigr) < \varepsilon \]
\emph{for all $x',x'' \in E$ with $0 < d_X(x',p) < \delta$ and
$0 < d_X(x'',p) < \delta$.}
\provenance{ECNU \cjk{定理}{Thm} 3.11 (bk.\ p.\,52); USTC \cjk{定理}{Thm} 9
(bk.\ pp.\,29--30). Rudin has no counterpart in Chapter 4.}

\begin{proof}[Proof sketch]
\emph{Necessity} is the triangle inequality: if $f(x) \to q$, take $\delta$
for $\varepsilon/2$ and add.

\emph{Sufficiency} routes through sequences, exactly as 4.2 does. Let
$(p_n)$ be any sequence in $E \setminus \{p\}$ with $p_n \to p$. The condition
makes $\bigl(f(p_n)\bigr)$ a Cauchy sequence, hence convergent by
completeness (3.11). If $(p_n')$ is a second such sequence with a different
limit, interleave the two into one sequence; its image is Cauchy, so all its
subsequential limits agree, and the two limits coincide. Now 4.2 applies.
\end{proof}

\begin{deeper}[Why Rudin can do without it, and you cannot]
Rudin never needs S4.1 because he almost never has to \emph{produce} a
function limit: his theorems take a convergent hypothesis and transport it.
The one place he does need it, he writes it out by hand --- Exercise 4.13,
extending a uniformly continuous function to the closure, is precisely
``Cauchy criterion + completeness'' done in a special case. Supplement E
below (S4.15) states the general fact the exercise is a shadow of.

Note what the criterion costs: completeness of the \emph{target}. Over $\Q$
it is false, and the failure is the same one that motivates all of
Chapter 1.
\end{deeper}

\sitem{S4.2}{Definition}
\emph{(Limits along a subset.) Let $E \subset X$, let $p$ be adherent to $E$
--- that is, $p \in \overline E$ --- and let $f$ be defined on $E$. Write}
\[ \lim_{x\to p,\ x\in E} f(x) = q \]
\emph{if for every $\varepsilon>0$ there is $\delta>0$ with
$d\bigl(f(x),q\bigr) < \varepsilon$ for all $x \in E$ satisfying $d(x,p) <
\delta$.}
\provenance{Tao I, Def.\ 9.3.6 (bk.\ p.\,221), with adherent points
Def.\ 9.1.8.}

\begin{unpack}[One definition in place of Rudin's four]
Tao's single notion specialises to everything Rudin defines separately.
Fix a function $f$ on a set $D$ and a point $p$:

\smallskip
\begin{tabular}{@{}ll@{}}
\toprule
\textbf{Take $E =$} & \textbf{and you get} \\
\midrule
$D \setminus \{p\}$        & Rudin's $\lim_{x\to p} f(x)$ (4.1) \\
$D \cap (p,\infty)$        & the right-hand limit $f(p+)$ (4.25) \\
$D \cap (-\infty,p)$       & the left-hand limit $f(p-)$ \\
$D$ itself                 & continuity at $p$, when the value is $f(p)$ \\
$D$, with $p = +\infty$    & $\lim_{x\to+\infty} f(x)$ (4.33) \\
$D \cap A$                 & the limit of the restriction $f|_A$ (4.12) \\
\bottomrule
\end{tabular}

\smallskip
The fourth row is the interesting one. Rudin's limit at $p$ is
\emph{deleted}: $p$ is thrown out of the domain, which is why 4.1 requires
$p$ to be a limit point and why $\lim_{x\to p}f(x)$ can differ from $f(p)$
--- the whole content of Definition 4.5 and of the discontinuity
classification. Tao's undeleted version with $E = D$ makes ``continuous at
$p$'' literally the statement ``the limit along the domain is $f(p)$.'' The
two conventions disagree, and each is chosen to make its own theorems
cleanest; knowing which one is in force is a standing source of confusion
between textbooks.
\end{unpack}

\sitem{S4.3}{Theorem}
\emph{(Squeeze.) If $f \le g \le h$ on $E$ and}
$\lim_{x\to p,\,x\in E} f = \lim_{x\to p,\,x\in E} h = L$,
\emph{then} $\lim_{x\to p,\,x\in E} g = L$.
\provenance{Tao I, Ex.\ 9.3.5; ECNU \cjk{定理}{Thm} 3.2 (bk.\ p.\,46).}
\begin{proof}
Let $(p_n)$ be any sequence in $E$ with $p_n \to p$. Then
$f(p_n) \to L$ and $h(p_n) \to L$, and $f(p_n) \le g(p_n) \le h(p_n)$, so
$g(p_n) \to L$ by the sequence squeeze (3.19 applied to
$g(p_n)-f(p_n) \le h(p_n)-f(p_n)$). Since the sequence was arbitrary, 4.2
gives the conclusion.
\end{proof}

\noindent
Rudin proves the sequence version (3.19 is its cousin) and never states the
function version, though he uses the reasoning freely.\kn{Cheap to prove and
constantly needed --- Supplement C below cannot get off the ground without
it. The mechanism is 4.2: pull back to sequences and use the sequence
squeeze.}

\sitem{S4.4}{Definition}
\emph{(Oscillation.) For $f$ real and bounded near $p$, put}
\[ \omega_f(p) = \lim_{r\to0^+}\ \Bigl[\ \sup_{d(x,p)<r} f(x)
   \;-\; \inf_{d(x,p)<r} f(x)\ \Bigr]. \]
\emph{The bracket decreases as $r$ decreases, so the limit exists.}
\provenance{Tao I, Ex.\ 9.3.4 and Tao II, Ex.\ 2.2.9 (as
$\limsup - \liminf$); Yu Pin, Lecture 21 Lemma 131 (p.\,224), who defines it
for maps of metric spaces as $\omega(f,p) = \inf_{U \ni p}
\operatorname{diam} f(U)$ over open $U$. USTC and ECNU introduce
\cjk{振幅}{zhenfu} only later, for integrability.}

\begin{center}
\begin{tikzpicture}[x=1mm, y=1mm]
  \draw[axis] (0,0) -- (52,0);
  \draw[seg] (4,5) .. controls (10,9) and (14,4) .. (18,7)
             .. controls (22,10) and (23,9) .. (25,12);
  \node[ptopen] at (25,12) {};
  \node[ptfill, minimum size=2.2pt] at (25,22) {};
  \draw[seg] (25,22) .. controls (28,25) and (31,21) .. (35,24)
             .. controls (41,28) and (44,25) .. (48,29);
  \draw[guide] (25,0) -- (25,31);
  \node[mlbl, anchor=north] at (25,-1.4) {$p$};
  \foreach \r/\yy in {12/33, 7/36.5, 3.5/40}{
    \draw[thin, {Stealth[length=2.5pt]}-{Stealth[length=2.5pt]}]
      (25-\r,\yy) -- (25+\r,\yy);}
  \node[mlbl, anchor=west] at (39,40) {shrinking windows};
  % bars
  \draw[axis] (60,0) -- (60,34);
  \foreach \x/\h/\lab in {66/26/{r_1}, 74/19/{r_2}, 82/15/{r_3}}{
    \draw[thin] (\x-2.4,0) rectangle (\x+2.4,\h);
    \node[mlbl, anchor=north] at (\x,-1.4) {$\lab$};}
  \draw[guide] (61,10) -- (90,10);
  \node[mlbl, anchor=west] at (90.6,10) {$\omega_f(p)$};
  \node[mlbl, anchor=south] at (74,35) {$\sup-\inf$ on each window};
\end{tikzpicture}
\end{center}
\figcap{Oscillation: measure $\sup - \inf$ over a window about $p$ and let
the window close. The bars fall but stall at the size of the jump. The
function is continuous at $p$ exactly when they fall to $0$.}

\sitem{S4.5}{Theorem}
\emph{$f$ is continuous at $p$ if and only if $\omega_f(p) = 0$. Moreover
$\Omega_\varepsilon = \{x : \omega_f(x) \ge \varepsilon\}$ is closed for each
$\varepsilon > 0$; consequently the set of points at which $f$ is
\emph{discontinuous} is $\bigcup_n \Omega_{1/n}$, a countable union of closed
sets, and the set where $f$ is continuous is a countable intersection of
open sets.}
\provenance{Yu Pin, Lemmas 131--132 and p.\,227, for arbitrary maps of metric
spaces; Tao I, Ex.\ 9.3.4.}
\begin{proof}
If $\omega_f(p)=0$ then given $\varepsilon>0$ some window has
$\sup - \inf < \varepsilon$, so $|f(x)-f(p)| < \varepsilon$ throughout it:
$f$ is continuous at $p$. Conversely continuity makes both $\sup$ and $\inf$
over a small enough window lie within $\varepsilon$ of $f(p)$, so their
difference is at most $2\varepsilon$.

For the second claim, let $\omega_f(x_0) < \varepsilon$. Then some ball
$B(x_0,r)$ has $\sup-\inf < \varepsilon$ on it, and every $x \in B(x_0,r)$
has a smaller ball inside $B(x_0,r)$, so $\omega_f(x) < \varepsilon$ too.
Hence $\{\omega_f < \varepsilon\}$ is open and its complement
$\Omega_\varepsilon$ is closed. Finally $f$ is discontinuous at $x$ exactly
when $\omega_f(x)>0$, that is when $x \in \Omega_{1/n}$ for some $n$.
\end{proof}

\begin{deeper}[The payoff: why Thomae's function has no mirror image]
Rudin's Exercise 4.18 produces a function continuous at every irrational and
discontinuous at every rational. The obvious next question --- can a function
be continuous at every \emph{rational} and discontinuous at every irrational?
--- Rudin never asks, and with his Chapter 4 tools he cannot answer it. S4.5
answers it: no.

The continuity set is always $\bigcap_n \{x : \omega_f(x) < 1/n\}$, an
intersection of countably many open sets. If some $f$ were continuous exactly
on $\Q$, then $\Q$ would be such an intersection, say $\Q = \bigcap_n U_n$
with each $U_n$ open and dense (each contains $\Q$). Enumerate
$\Q = \{q_1,q_2,\dots\}$; then
\[ \varnothing = \Bigl(\bigcap_n U_n\Bigr) \cap
   \Bigl(\bigcap_k \bigl(\R\setminus\{q_k\}\bigr)\Bigr) \]
exhibits $\varnothing$ as a countable intersection of dense open sets, which
Baire's theorem forbids. So Thomae's function has no mirror image, and the
asymmetry between $\Q$ and its complement --- the theme running from 2.44 to
Chapter 3's rearrangements --- shows up once more.

Rudin proves Baire's theorem in Exercise 3.22; this is one of its first real
uses.
\end{deeper}

% ============================================================
\ssection{Supplement B --- Orders of magnitude: the calculus of limits}

\begin{roadmap}[Why this supplement belongs to Chapter 4]
\raggedright
\textbf{What Rudin does.} Nothing --- and that is the point. There is no
$o$, no $O$, no notion of one quantity being negligible beside another
anywhere in Chapters 1--4. Limits are proved to exist and are then named.

\smallskip
\textbf{The gap.} Analysis outside this book is conducted almost entirely in
asymptotic notation, and the notation is not decoration: it is how one keeps
track of which error terms matter. Rudin's abstract line never needs it
because he never computes a limit that is not already a theorem.

\smallskip
\textbf{Where the gap bites.} At 5.15, where Taylor's theorem is really a
statement that a remainder is $o\bigl((x-a)^n\bigr)$ and has to be phrased
instead with an explicit intermediate point; at 3.20, where the special
limits are exactly statements about relative rates of growth; and in every
computation a reader will meet outside these covers.

\smallskip
\textbf{What this section supplies.} The definitions, the substitution
theorem that makes them usable, and --- most importantly --- the pitfall that
substitution is legal in factors and illegal in summands, which is the error
the notation invites.
\end{roadmap}

\sitem{S4.6}{Definition}
\emph{Let $f,g$ be defined near $p$ (with $g \ne 0$ there). As $x \to p$:}
\begin{enumerate}[label=(\alph*), leftmargin=2.2em, itemsep=2pt, topsep=3pt]
  \item $f = o(g)$ \emph{means} $f/g \to 0$;
  \item $f = O(g)$ \emph{means} $|f| \le C|g|$ \emph{near $p$, for some
        constant $C$};
  \item $f \sim g$ (\emph{equivalent}) \emph{means} $f/g \to 1$.
\end{enumerate}
\provenance{ECNU \cjk{无穷小量的阶}{orders of infinitesimals}, bk.\ pp.\,57--59;
USTC \cjk{等价无穷小替换}{equivalent-infinitesimal substitution}, bk.\
pp.\,32--35.}

\noindent
The equality sign in $f = o(g)$ is an abuse: $o(g)$ names a \emph{class} of
functions and the statement is membership.\kn{Read every ``$=$'' involving
$o$ or $O$ left-to-right and never reverse it. From $f = o(x)$ and
$h = o(x)$ it does not follow that $f = h$; both are merely in the same
class. ECNU makes this explicit by writing $o(g) = \{f : f/g \to 0\}$.}

\sitem{S4.7}{Theorem}
\emph{(Substitution of equivalents.) Suppose $\alpha \sim \beta$ as
$x \to p$. If $\lim \alpha u$ exists, then $\lim \beta u$ exists and they are
equal; if $\lim v/\alpha$ exists, then $\lim v/\beta$ exists and they are
equal.}
\provenance{ECNU \cjk{定理}{Thm} 3.12; USTC \cjk{定理}{Thm} 12.}
\begin{proof}
Write $\beta = \alpha \cdot (\beta/\alpha)$ where $\beta/\alpha \to 1$. Then
$\beta u = (\alpha u)(\beta/\alpha)$, and the limit of a product is the
product of the limits (4.4), so $\lim \beta u = (\lim \alpha u)\cdot 1$.
The quotient case is identical with $v/\beta = (v/\alpha)(\alpha/\beta)$.
\end{proof}

\begin{pitfall}[Substitution is legal in products and quotients --- nowhere else]
Since $\tan x \sim x$ and $\sin x \sim x$ as $x \to 0$, one is tempted to
replace both in a \emph{difference}:
\[ \tan x - \sin x \;\overset{?}{\sim}\; x - x = 0 , \]
which is not merely imprecise but wrong. The truth is
\[ \tan x - \sin x = \sin x\,\frac{1-\cos x}{\cos x}
   \sim x \cdot \frac{x^2}{2} = \frac{x^3}{2}, \]
so $(\tan x - \sin x)/\sin^3 x \to \tfrac12$, not $0$. The two equivalences
are each accurate to first order, and the difference cancels precisely the
order they are accurate to; what survives is third order, about which they
say nothing.

The rule to carry: equivalents may be substituted into a \emph{factor} of a
product or quotient, never into a summand. When a difference is what you
want, expand to enough orders --- which is what Taylor's theorem (5.15) will
industrialise.
\end{pitfall}

\sitem{S4.8}{Definition}
\emph{$f$ is an \emph{infinitely large quantity} as $x \to p$ if
$|f(x)| \to +\infty$.}\kn{Not the same as unbounded, and the difference is
the point of the definition. On $(0,\infty)$ the function $x\sin x$ is
unbounded, yet it does not tend to $\infty$: at $x = n\pi$ it is $0$. ECNU
isolates this as the standard beginner's error --- Rudin's 4.33 defines the
notion correctly but never contrasts it with unboundedness.}

\sitem{S4.9}{Definition}
\emph{(Asymptotes.) A line $L$ is an \emph{asymptote} of the curve $y=f(x)$
if the distance from $\bigl(x,f(x)\bigr)$ to $L$ tends to $0$ as the point
recedes to infinity. The line $y = ax+b$ is an asymptote as $x \to +\infty$
if and only if}
\[ a = \lim_{x\to+\infty}\frac{f(x)}{x},
   \qquad b = \lim_{x\to+\infty}\bigl[f(x)-ax\bigr], \]
\emph{both limits existing finitely; and $x = c$ is a vertical asymptote if
some one-sided limit of $f$ at $c$ is infinite.}
\provenance{ECNU \cjk{曲线的渐近线}{asymptotes}, bk.\ pp.\,61--62; USTC
Ex.\ 1.2.14. Absent from Rudin entirely.}

\begin{center}
\begin{tikzpicture}[x=1mm, y=1mm]
  \draw[axis] (0,0) -- (96,0);
  \draw[axis] (0,-3) -- (0,40);
  \draw[guide] (10,0) -- (10,38);
  \node[mlbl, anchor=north] at (10,-1.4) {$x=c$};
  \draw[thin] (20,0) -- (92,28.8);
  \node[mlbl, anchor=north west] at (78,22.4) {$y=ax+b$};
  \draw[seg] (18,34.2) .. controls (21,27) and (23,22) .. (30,18)
             .. controls (36,15.6) and (40,16.4) .. (45,18)
             .. controls (52,20) and (60,22.6) .. (70,24.7)
             .. controls (78,26.4) and (84,28) .. (90,29.6);
  \draw[thin, {Stealth[length=2.5pt]}-{Stealth[length=2.5pt]}]
    (60,16) -- (60,21.9);
  \node[mlbl, anchor=west] at (62,13.6) {$f(x)-(ax+b)\to0$};
\end{tikzpicture}
\end{center}
\figcap{An oblique asymptote. Because the distance from $(x,f(x))$ to the
line is $|f(x)-ax-b|/\sqrt{1+a^2}$, the geometric condition ``distance
$\to 0$'' is the algebraic condition $f(x)-ax-b \to 0$ --- which is why
dividing by $x$ recovers $a$ first, and subtracting $ax$ then recovers $b$.
A vertical asymptote is the separate case where $f$ itself blows up.}

\begin{unpack}[Worked example]
For $f(x) = \dfrac{x^3}{x^2+2x-3}$ division gives
\[ f(x) = x-2 + \frac{7x-6}{x^2+2x-3}, \]
and the remainder tends to $0$, so $y = x-2$ is an asymptote in both
directions. Since $x^2+2x-3 = (x+3)(x-1)$, the lines $x=-3$ and $x=1$ are
vertical asymptotes. Note that the recipe found $a$ and $b$ without any
division: $f(x)/x \to 1$ and $f(x)-x \to -2$.
\end{unpack}

% ============================================================
\ssection{Supplement C --- The two limits Rudin postpones}

\begin{pitfall}[Read this section as a loan against Chapter 8]
Rudin has no $\sin$, $\cos$, $\exp$ or $\log$ at this point: they are
\emph{defined} in 8.7 by power series, and every property of them is derived
there. So the two theorems below are not available in his development, and
the geometric proof of the first uses the area of a circular sector --- that
is, integration (Chapter 6) and arc length (8.7 again).

They are included because every other text proves them here, because every
computation in a first course needs them, and because the arguments are
worth seeing. But the logical order is Rudin's, and Yu Pin --- who also
defines $\sin$ by its series --- shows what that costs: he proves the same
limit in one line from the estimate
\[ \Bigl|\frac{\sin x}{x} - 1\Bigr| \le |x|^2 e^{|x|}, \]
read straight off the series with no geometry at all. The picture below is
then an illustration of a fact already proved, not the proof.
\end{pitfall}

\sitem{S4.10}{Theorem}
\emph{$\displaystyle\lim_{x\to0}\frac{\sin x}{x} = 1$.}
\provenance{ECNU \cjk{定理}{Thm} 3.14 (bk.\ p.\,53); USTC \cjk{定理}{Thm} 10
(bk.\ p.\,30). Rudin obtains it from 8.7.}

\begin{proof}[The geometric proof]
Let $0 < x < \pi/2$ and work in the unit circle. Comparing the areas of the
triangle $OAB$, the sector $OAB$, and the triangle $OAT$ gives
\[ \tfrac12\sin x \;<\; \tfrac12 x \;<\; \tfrac12\tan x . \]
Dividing by $\tfrac12\sin x > 0$ yields $1 < x/\sin x < 1/\cos x$, that is
\[ \cos x < \frac{\sin x}{x} < 1 . \]
Since $\cos x \to 1$ as $x \to 0^+$, the squeeze S4.3 gives the limit from
the right; the quotient is even, so the two-sided limit follows.
\end{proof}

\begin{center}
\begin{tikzpicture}[x=1mm, y=1mm]
  \fill[black!26] (0,0) -- (34,0)
    arc[start angle=0, end angle=55, radius=34] -- cycle;
  \fill[black!10] (0,0) -- (34,0) -- (19.50,27.85) -- cycle;
  \fill[pattern=north east lines, pattern color=black!45]
    (34,0) -- (34,48.56) -- (19.50,27.85)
    arc[start angle=55, end angle=0, radius=34] -- cycle;
  \draw[thin] (0,0) -- (34,0);
  \draw[thin] (0,0) -- (34,48.56);
  \draw[thin] (34,0) -- (34,48.56);
  \draw[thin] (19.50,27.85) -- (34,0);
  \draw[seg] (34,0) arc[start angle=0, end angle=55, radius=34];
  \node[ptfill, minimum size=2.2pt] at (19.50,27.85) {};
  \node[mlbl, anchor=north east] at (-0.6,0.4) {$O$};
  \node[mlbl, anchor=north] at (34,-1.6) {$A$};
  \node[mlbl, anchor=south east] at (18.6,28.4) {$B$};
  \node[mlbl, anchor=south] at (34,50.2) {$T$};
  \node[mlbl] at (11,3.6) {$x$};
  % legend
  \fill[black!10] (46,42) rectangle (52,46.4);
  \draw[thin] (46,42) rectangle (52,46.4);
  \node[lbl, anchor=west] at (53.4,44.2)
    {triangle $OAB$: area $\tfrac12\sin x$};
  \fill[black!26] (46,32) rectangle (52,36.4);
  \draw[thin] (46,32) rectangle (52,36.4);
  \node[lbl, anchor=west] at (53.4,34.2)
    {${}+{}$ sliver ${}={}$ sector: area $\tfrac12 x$};
  \fill[pattern=north east lines, pattern color=black!45]
    (46,22) rectangle (52,26.4);
  \draw[thin] (46,22) rectangle (52,26.4);
  \node[lbl, anchor=west] at (53.4,24.2)
    {${}+{}$ this ${}={}$ triangle $OAT$: area $\tfrac12\tan x$};
\end{tikzpicture}
\end{center}
\figcap{Three nested regions in the unit circle. Containment of the regions
\emph{is} the chain $\sin x < x < \tan x$ --- the only inequality the proof
needs. The sector's area being $\tfrac12x$ is what makes radian measure the
right measure, and is also what this argument borrows from Chapter 6.}

\sitem{S4.11}{Theorem}
\emph{$\displaystyle\lim_{x\to\infty}\Bigl(1+\frac1x\Bigr)^{x} = e$, in both
directions; equivalently $\lim_{t\to0}(1+t)^{1/t} = e$.}
\provenance{ECNU bk.\ p.\,54; USTC \cjk{定理}{Thm} 11. Rudin proves only the
sequence case, 3.31.}
\begin{proof}
Let $x>0$ and $n = \lfloor x \rfloor$, so $n \le x < n+1$. Increasing the
base increases the power, and increasing the exponent does too, so
\[ \Bigl(1+\frac{1}{n+1}\Bigr)^{n}
   \;<\; \Bigl(1+\frac1x\Bigr)^{x}
   \;<\; \Bigl(1+\frac1n\Bigr)^{n+1} . \]
The left side is $\bigl(1+\tfrac1{n+1}\bigr)^{n+1}$ divided by
$\bigl(1+\tfrac1{n+1}\bigr)$ and the right side is
$\bigl(1+\tfrac1n\bigr)^{n}$ times $\bigl(1+\tfrac1n\bigr)$; by 3.31 both
tend to $e$. Since $n \to \infty$ as $x \to +\infty$, the squeeze S4.3 gives
the limit. For $x \to -\infty$ put $x = -(y+1)$ and compute
$\bigl(1+\tfrac1x\bigr)^{x} = \bigl(1+\tfrac1y\bigr)^{y+1} \to e$.
\end{proof}

\begin{unpack}[From Rudin's sequence to a continuous variable]
Rudin's 3.31 establishes $\bigl(1+\tfrac1n\bigr)^n \to e$ for integers $n$.
To pass to a real variable, trap $x$ between consecutive integers: with
$n = \lfloor x \rfloor$ and $x>0$,
\[ \Bigl(1+\frac{1}{n+1}\Bigr)^{n} \;<\;
   \Bigl(1+\frac{1}{x}\Bigr)^{x} \;<\;
   \Bigl(1+\frac{1}{n}\Bigr)^{n+1}, \]
because increasing the base increases the power and so does increasing the
exponent. Both outer terms are the sequence limit multiplied by a factor
tending to $1$ --- the left is $\bigl(1+\tfrac1{n+1}\bigr)^{n+1}$ divided by
$\bigl(1+\tfrac1{n+1}\bigr)$, the right is $\bigl(1+\tfrac1n\bigr)^n$ times
$\bigl(1+\tfrac1n\bigr)$ --- so both tend to $e$, and S4.3 closes the gap.
The case $x \to -\infty$ follows from the substitution $x = -(y+1)$.

This bracketing is the standard way to lift \emph{any} sequence limit to a
continuous one when the function is monotone between integers, and it is
worth keeping for that reason alone.
\end{unpack}

% ============================================================
\ssection{Supplement D --- Monotonicity, inverses, and the elementary functions}

\begin{roadmap}[Why this supplement belongs to Chapter 4]
\raggedright
\textbf{What Rudin does.} Theorem 4.17 proves that a continuous bijection of
a \emph{compact} metric space has a continuous inverse, and Example 4.21
shows compactness cannot simply be dropped. Monotone functions appear
separately at 4.28--4.30, and the elementary functions do not appear at all
--- $\exp$, $\log$, $\sin$ are built in Chapter 8 from power series.

\smallskip
\textbf{The gap.} On the line there is a second route to the same conclusion
that uses no compactness whatever, only order --- and it reaches domains
4.17 cannot, such as open and unbounded intervals. Rudin's two topics,
inverse functions and monotone functions, are in fact one topic, and he never
connects them.

\smallskip
\textbf{Where the gap bites.} At 4.21 itself, where the circle-wrap
counterexample is offered as evidence that compactness is \emph{needed},
when what it really shows is that the domain was not ordered; and at
Exercise 4.14, which is about monotone open maps and sits next to a theorem
that would explain it.

\smallskip
\textbf{What this section supplies.} The equivalence of injectivity and
strict monotonicity for continuous functions on an interval; the inverse
function theorem that follows, with no compactness; and the catalogue of
elementary functions that this route makes available four chapters early.
\end{roadmap}

\sitem{S4.12}{Theorem}
\emph{If $f$ is continuous and one-to-one on an interval $I \subset \R$, then
$f$ is strictly monotone on $I$.}
\provenance{Tao I, Ex.\ 9.8.3; USTC \cjk{定理}{Thm} 3 states the equivalence
(continuous $f$ on $[a,b]$ is invertible iff strictly monotone).}
\begin{proof}
Suppose $f$ is not strictly monotone. Then there are $u<v<w$ in $I$ with
$f(v)$ not strictly between $f(u)$ and $f(w)$ --- say $f(v) \ge f(u)$ and
$f(v) \ge f(w)$, the other case being symmetric. Pick $c$ with
$\max\{f(u),f(w)\} < c < f(v)$, possible unless $f(v)$ equals one of them,
which injectivity forbids. The intermediate value theorem 4.23 applied on
$[u,v]$ and on $[v,w]$ produces $s \in (u,v)$ and $t \in (v,w)$ with
$f(s)=f(t)=c$, contradicting injectivity.
\end{proof}

\noindent
The proof is the intermediate value theorem used three times: if $f$ were
neither increasing nor decreasing, some triple $u<v<w$ would have $f(v)$
outside the interval with endpoints $f(u)$, $f(w)$, and 4.23 would then
produce a repeated value on one side or the other.\kn{The converse is
trivial, so this says: for continuous functions on an interval, injectivity
and strict monotonicity are the \emph{same} hypothesis. Rudin's Exercise 4.14
(open maps are monotone) is a relative of this; the equivalence itself he
never states.}

\sitem{S4.13}{Theorem}
\emph{(Inverse functions without compactness.) Let $f$ be continuous and
strictly increasing on an interval $I$. Then $J = f(I)$ is an interval, $f :
I \to J$ is a bijection, and $f^{-1} : J \to I$ is continuous and strictly
increasing.}
\provenance{ECNU \cjk{定理}{Thm} 4.8 (bk.\ pp.\,73--74); USTC \cjk{定理}{Thm}
3; Tao I, Prop.\ 9.8.3.}

\begin{proof}[The $\delta$ that monotonicity hands you]
Let $y_0 = f(x_0)$ with $x_0$ interior to $I$, and let $\varepsilon > 0$ be
small enough that $[x_0-\varepsilon, x_0+\varepsilon] \subset I$. Put
$y_1 = f(x_0-\varepsilon)$ and $y_2 = f(x_0+\varepsilon)$; strict increase
gives $y_1 < y_0 < y_2$. Take
\[ \delta = \min\{\,y_0-y_1,\ y_2-y_0\,\} > 0 . \]
If $|y - y_0| < \delta$ then $y_1 < y < y_2$, and applying the increasing
map $f^{-1}$ gives $x_0-\varepsilon < f^{-1}(y) < x_0+\varepsilon$. Endpoints
of $I$ are handled one-sidedly.
\end{proof}

\begin{center}
\begin{tikzpicture}[x=1mm, y=1mm]
  \draw[axis] (0,0) -- (76,0);
  \draw[axis] (0,0) -- (0,48);
  \draw[seg] (6,4) .. controls (16,10) and (20,13) .. (30,20)
             .. controls (40,27) and (48,33) .. (58,40)
             .. controls (63,43) and (66,45) .. (70,46.5);
  \draw[guide] (24,0) -- (24,16) -- (0,16);
  \draw[guide] (34,0) -- (34,24) -- (0,24);
  \draw[guide] (44,0) -- (44,31) -- (0,31);
  \node[mlbl, anchor=north] at (24,-1.4) {$x_0{-}\varepsilon$};
  \node[mlbl, anchor=north] at (34,-1.4) {$x_0$};
  \node[mlbl, anchor=north] at (44,-1.4) {$x_0{+}\varepsilon$};
  \node[mlbl, anchor=east] at (-1.4,16) {$y_1$};
  \node[mlbl, anchor=east] at (-1.4,24) {$y_0$};
  \node[mlbl, anchor=east] at (-1.4,31) {$y_2$};
  \draw[thin, {Stealth[length=2.5pt]}-{Stealth[length=2.5pt]}]
    (68,17) -- (68,31);
  \node[mlbl, anchor=west] at (69.4,24) {$\delta$};
  \draw[guide] (0,17) -- (68,17);
  \draw[guide] (44,31) -- (68,31);
\end{tikzpicture}
\end{center}
\figcap{Theorem S4.13: the $\varepsilon$-interval about $x_0$ is carried by
$f$ onto an interval about $y_0$, and $\delta$ is simply the smaller of its
two halves. Monotonicity is doing all the work --- it is what makes the
image of an interval an interval, and what lets $f^{-1}$ preserve the
trapping. No compactness enters, so $I$ may be open or unbounded.}

\begin{deeper}[Two theorems about inverses, neither containing the other]
Rudin 4.17 and S4.13 are genuinely different results:

\smallskip
\begin{tabular}{@{}p{0.30\linewidth}p{0.30\linewidth}p{0.30\linewidth}@{}}
\toprule
& \textbf{Rudin 4.17} & \textbf{S4.13} \\
\midrule
domain      & compact metric space   & interval in $\R$ \\
hypothesis  & continuous bijection   & continuous, strictly monotone \\
mechanism   & compactness (4.14, 2.35)& order + IVT \\
reaches     & $\R^k$, function spaces & $(0,\infty)$, $\R$, half-open \\
misses      & $(0,1) \to (0,1)$       & anything not ordered \\
\bottomrule
\end{tabular}

\smallskip
Example 4.21 --- the circle-wrap $t \mapsto (\cos t, \sin t)$ on $[0,2\pi)$,
whose inverse is discontinuous --- shows 4.17 fails without compactness; but
$f(x) = x/(1-x^2)$ on $(-1,1)$ is a continuous increasing bijection onto
$\R$ with continuous inverse, on a domain that is not compact. Order
supplies what compactness supplied.
\end{deeper}

\sitem{S4.14}{Theorem}
\emph{(Continuity of the elementary functions.) For $a>0$ the function
$a^x$, defined for irrational $x$ as $\sup\{a^r : r \in \Q,\ r \le x\}$, is
continuous on $\R$ and satisfies $a^{\alpha}a^{\beta}=a^{\alpha+\beta}$ and
$(a^{\alpha})^{\beta}=a^{\alpha\beta}$. Its inverse $\log_a$ is continuous
on $(0,\infty)$ by S4.13, the power $x^{\alpha}=e^{\alpha\ln x}$ is
continuous on $(0,\infty)$, and $\sin$, $\cos$ are continuous because
$|\sin x - \sin y| \le |x-y|$. Consequently every function built from these
by finitely many algebraic operations, compositions, and inversions is
continuous on its domain of definition.}
\provenance{ECNU \cjk{定理}{Thm} 4.10--4.13 (bk.\ pp.\,78--80); USTC
\cjk{定理}{Thm} 4 (bk.\ pp.\,45--46).}
\begin{proof}[Proof sketch]
Continuity of $a^x$ at $0$ follows from $a^{1/n}\to1$ (Rudin 3.20(b)) and
monotonicity; the identity $a^{x+h}=a^xa^h$ then transports continuity from
$0$ to every point, since $|a^{x+h}-a^x| = a^x|a^h-1|$. The laws
$a^{\alpha}a^{\beta}=a^{\alpha+\beta}$ and
$(a^{\alpha})^{\beta}=a^{\alpha\beta}$ hold for rational exponents by S1.13
and extend to real exponents by density and continuity (Rudin's Exercise
4.4). Continuity of $\log_a$ is S4.13, that of $x^{\alpha}=e^{\alpha\ln x}$
is composition, and $|\sin x - \sin y| \le |x-y|$ gives the trigonometric
cases. Closure under sums, products, quotients and composition is 4.9 and
4.7.
\end{proof}

\begin{deeper}[The same destination by two routes]
Rudin builds these functions in Chapter 8 from power series, and their
continuity is then a corollary of the theorem that a power series is
continuous inside its disc of convergence (8.1). The Chinese texts build
them here, from the supremum definition of $a^x$ that Rudin sets as
Exercises 1.6--1.7, and get continuity by hand.

The trade is visible: Rudin's route is shorter and yields differentiability
and the addition formulas at the same stroke, but it defers every concrete
function by four chapters; the elementary route has functions available
immediately but must prove each property separately. Neither is more
rigorous. The Chinese texts' order is why their students compute limits in
week three and Rudin's readers do not.
\end{deeper}

% ============================================================
\ssection{Supplement E --- Uniform continuity: the working toolkit}

\begin{roadmap}[Why this supplement belongs to Chapter 4]
\raggedright
\textbf{What Rudin does.} Definition 4.18 gives uniform continuity, 4.19
proves that continuity on a compact set implies it, and 4.20 exhibits the
standard failures. Everything else about the notion is set as Exercises
4.8--4.13.

\smallskip
\textbf{The gap.} Those exercises contain the working theory --- uniform
continuity preserves Cauchy sequences, preserves boundedness, permits
extension to the closure --- and as exercises they are unavailable to a
reader who has not done them, and unquotable to one who has. There is also
no \emph{criterion}: to show a function is not uniformly continuous, Rudin
negates the definition by hand each time.

\smallskip
\textbf{Where the gap bites.} In Chapter 8, where $b^x$ is defined on $\Q$
and must be extended to $\R$ --- an appeal to Exercise 4.13; and at 7.16 and
beyond, where uniform continuity is used as a known tool.

\smallskip
\textbf{What this section supplies.} The sequential criterion, which turns
every non-example into a two-line computation; the gluing lemma; the
extension theorem in its general form; the algebra of what preserves uniform
continuity; and the Lebesgue number lemma, which is what Rudin's own proof of
4.19 manufactures without naming.
\end{roadmap}

\sitem{S4.15}{Theorem}
\emph{(Sequential criterion.) $f : E \to Y$ is uniformly continuous on $E$
if and only if for every pair of sequences $(p_n)$, $(q_n)$ in $E$ with
$d(p_n,q_n) \to 0$ one has $d\bigl(f(p_n),f(q_n)\bigr) \to 0$.}
\provenance{Tao I, Prop.\ 9.9.8 (via ``equivalent sequences''); ECNU
\cjk{例}{Ex} 10 (bk.\ p.\,75).}
\begin{proof}
If $f$ is uniformly continuous and $d(p_n,q_n)\to0$, then given
$\varepsilon>0$ take $\delta$ from the definition; eventually
$d(p_n,q_n)<\delta$, so eventually
$d\bigl(f(p_n),f(q_n)\bigr)<\varepsilon$.

Conversely, suppose $f$ is not uniformly continuous. Then some
$\varepsilon>0$ resists every $\delta$; taking $\delta = 1/n$ produces
points $p_n,q_n$ with $d(p_n,q_n)<1/n$ but
$d\bigl(f(p_n),f(q_n)\bigr) \ge \varepsilon$. This pair of sequences
violates the stated condition.
\end{proof}

\noindent
The value is in the negation: to prove a function \emph{not} uniformly
continuous, exhibit one pair of sequences that close up while their images do
not.\kn{Compare Rudin's own method at 4.20: he argues directly from the
$\varepsilon$-$\delta$ negation each time. The criterion packages that
argument once. Note the parallel it draws: continuity preserves
\emph{convergent} sequences (4.2), uniform continuity preserves
\emph{pairs that close up}.}

\begin{center}
\begin{tikzpicture}[x=1mm, y=1mm]
  \draw[axis] (0,0) -- (68,0);
  \draw[axis] (0,0) -- (0,40);
  \draw[seg] (5.5,36) .. controls (7,26) and (9,20) .. (13,15)
             .. controls (18,10.5) and (24,8) .. (32,6.4)
             .. controls (42,4.8) and (54,3.8) .. (64,3.4);
  \draw[guide] (6,0) -- (6,33) -- (0,33);
  \draw[guide] (12,0) -- (12,16.5) -- (0,16.5);
  \node[mlbl, anchor=north] at (6,-1.4) {$q_n$};
  \node[mlbl, anchor=north] at (13.6,-1.4) {$p_n$};
  \node[mlbl, anchor=east] at (-1.4,33) {$2n$};
  \node[mlbl, anchor=east] at (-1.4,16.5) {$n$};
  \draw[thin, {Stealth[length=2.5pt]}-{Stealth[length=2.5pt]}]
    (6,-6) -- (12,-6);
  \node[mlbl, anchor=north] at (9,-7.4) {$\to0$};
  \draw[thin, {Stealth[length=2.5pt]}-{Stealth[length=2.5pt]}]
    (-9,16.5) -- (-9,33);
  \node[mlbl, anchor=south, rotate=90] at (-10.4,25) {$\to\infty$};
  \node[lbl, anchor=west] at (34,22) {$f(x)=1/x$ on $(0,1]$};
  \node[lbl, anchor=west] at (34,17.6) {$p_n=1/n$, $q_n=1/2n$};
\end{tikzpicture}
\end{center}
\figcap{The criterion applied to $1/x$: the two arguments close up
($p_n - q_n \to 0$) while their values pull apart ($n$ versus $2n$). One
pair of sequences settles it --- no $\varepsilon$-$\delta$ negation needed.
For $x^2$ on $\R$ take $p_n = n$, $q_n = n + 1/n$: the gap tends to $0$ and
the images differ by $2 + 1/n^2$.}

\sitem{S4.16}{Theorem}
\emph{(Gluing.) If $a < c < b$ and $f$ is uniformly continuous on $[a,c]$
and on $[c,b]$, then $f$ is uniformly continuous on $[a,b]$. The same holds
for two intervals sharing one endpoint, bounded or not.}
\provenance{ECNU \cjk{例}{Ex} 12 (bk.\ p.\,76).}
\begin{proof}
Given $\varepsilon>0$, take $\delta_1$ and $\delta_2$ for $\varepsilon/2$ on
$[a,c]$ and on $[c,b]$, and put $\delta = \min\{\delta_1,\delta_2\}$. Let
$|x-y|<\delta$. If $x,y$ lie in the same piece there is nothing to prove.
Otherwise $x \le c \le y$, so $|x-c|<\delta$ and $|c-y|<\delta$, and
\[ |f(x)-f(y)| \le |f(x)-f(c)| + |f(c)-f(y)|
   < \tfrac{\varepsilon}{2}+\tfrac{\varepsilon}{2} = \varepsilon . \]
\end{proof}

\begin{unpack}[Why gluing needs an argument at all, and what it buys]
Uniform continuity is not a local property, so the union of two sets on which
$f$ is uniformly continuous is not automatically another. The proof works
because the two pieces meet: given $\varepsilon$, take $\delta_1,\delta_2$
for $\varepsilon/2$ on the two pieces and set
$\delta = \min\{\delta_1,\delta_2\}$. If $|x-y| < \delta$ with $x \le c \le
y$, then $|x-c| < \delta$ and $|c-y| < \delta$, so
\[ |f(x)-f(y)| \le |f(x)-f(c)| + |f(c)-f(y)| < \varepsilon . \]
The shared point $c$ is the whole mechanism, and it is why two intervals
with a \emph{gap} between them cannot be glued.

The payoff is a standard trick: to prove $f$ uniformly continuous on
$[0,\infty)$, prove it on the compact $[0,1]$ by 4.19 and on $[1,\infty)$ by
a direct estimate. For $f(x)=\sqrt x$ the second is
$|\sqrt x - \sqrt y| = |x-y|/(\sqrt x+\sqrt y) \le \tfrac12|x-y|$. Rudin's
Exercise 4.8 asks for the boundedness half of this story and never supplies
the gluing tool.
\end{unpack}

\sitem{S4.17}{Theorem}
\emph{(Extension to the closure.) Let $f : E \to Y$ be uniformly continuous
with $Y$ complete. Then $\lim_{x\to p} f(x)$ exists at every $p \in
\overline E$, and $f$ extends uniquely to a uniformly continuous function on
$\overline E$.}
\provenance{Tao I, Cor.\ 9.9.14 and Prop.\ 9.9.12. Rudin poses the special
case as Exercise 4.13.}
\begin{proof}
Let $p \in \overline E$ and let $(p_n)$ be a sequence in $E$ with
$p_n \to p$. It is Cauchy, and a uniformly continuous map sends Cauchy
sequences to Cauchy sequences: given $\varepsilon$, take $\delta$ from
uniform continuity and then $N$ with $d(p_n,p_m)<\delta$ for $n,m\ge N$.
Completeness of $Y$ makes $\bigl(f(p_n)\bigr)$ converge. Interleaving two
such sequences shows the limit does not depend on the choice, so
$\lim_{x\to p}f(x)$ exists by 4.2.

Defining $\tilde f(p)$ to be that limit extends $f$ to $\overline E$; the
extension is uniformly continuous because the defining inequality passes to
limits, and it is unique because two continuous maps agreeing on the dense
set $E$ agree on $\overline E$ (Rudin's Exercise 4.4).
\end{proof}

\begin{proofwalk}[How the proof flows]
  \item \textbf{Choose the method.} A \emph{value} has to be produced at a
        point where $f$ is not defined, so something must manufacture it.
        The only manufacturing device available is completeness of the
        target, and the only way to reach it is through Cauchy sequences.
  \item \textbf{Convert the hypothesis into a Cauchy statement.} Uniform
        continuity says one $\delta$ works everywhere, which is exactly what
        lets a Cauchy sequence in the domain be pushed to a Cauchy sequence
        in the target. Ordinary continuity cannot do this --- the $\delta$
        would move with the point.
  \item \textbf{Spend completeness.} Cauchy becomes convergent. This is the
        only step that uses anything about $Y$.
  \item \textbf{Remove the dependence on the sequence.} Interleaving two
        approximating sequences is the standard device; it converts ``every
        sequence gives a limit'' into ``the limit exists,'' which is 4.2.
  \item \textbf{Audit.} Uniform continuity is spent once, in step 2;
        completeness once, in step 3; density of $E$ in $\overline E$ once,
        to know the approximating sequences exist at all. Drop uniformity and
        the theorem fails --- $\sin(1/x)$ on $(0,1)$ is continuous, sends a
        Cauchy sequence to a divergent one, and extends to no value at $0$.
\end{proofwalk}

\noindent
This is the theorem Rudin's Exercise 4.13 is a special case of, and the
reason uniform continuity matters at all in Chapter 8: it is what lets a
function defined on $\Q$ --- like $b^x$ for rational $x$ in Exercise 1.6 ---
become a function on $\R$.\kn{Chain of reasoning: uniformly continuous maps
send Cauchy sequences to Cauchy sequences (Rudin's Exercise 4.11), the
target is complete, so the images converge; S4.1 then upgrades ``every
sequence works'' to ``the limit exists.'' Plain continuity is not enough:
$1/x$ on $(0,1)$ is continuous, sends the Cauchy sequence $(1/n)$ to an
unbounded one, and does not extend to $0$.}

\sitem{S4.18}{Theorem}
\emph{(What preserves uniform continuity.) If $f,g$ are uniformly continuous
on $E$ then so are $f+g$, $f-g$, $cf$, $\max(f,g)$, $\min(f,g)$, and $g\circ
f$ where defined. The product $fg$ need not be, but is if $f$ and $g$ are
both bounded.}
\provenance{Tao I, Ex.\ 9.9.6; Tao II, Ex.\ 2.3.4--2.3.6.}
\begin{proof}
For $f+g$ take $\delta$ to be the smaller of the two $\delta$'s for
$\varepsilon/2$; for $cf$ use $\varepsilon/|c|$; for $g \circ f$ feed the
$\delta$ of $g$ into $f$ as its $\varepsilon$. The identity
$\max(f,g) = \tfrac12(f+g+|f-g|)$ and $\bigl||u|-|v|\bigr| \le |u-v|$ handle
$\max$ and $\min$.

For products with both factors bounded by $M$,
\begin{align*}
  |f(p)g(p)-f(q)g(q)|
    &\le |f(p)|\,|g(p)-g(q)| + |g(q)|\,|f(p)-f(q)|\\
    &\le M\bigl(|g(p)-g(q)|+|f(p)-f(q)|\bigr),
\end{align*}
so $\delta$ for $\varepsilon/2M$ in each factor works. Without boundedness
the conclusion fails: $f(x)=g(x)=x$ are uniformly continuous on $\R^1$ and
$fg=x^2$ is not.
\end{proof}

\noindent
The counterexample is the smallest one available: $f(x)=g(x)=x$ is uniformly
continuous on $\R$, and $fg = x^2$ is not --- the same function Rudin uses at
4.20 for a different purpose.\kn{The boundedness repair is the usual
splitting: $|fg(p)-fg(q)| \le |f(p)|\,|g(p)-g(q)| + |g(q)|\,|f(p)-f(q)|$,
and if both factors are bounded by $M$ each term is small. Compare 3.3(c),
where exactly this decomposition proves the product rule for limits ---
the strips-and-corner picture there is the same algebra.}

\sitem{S4.19}{Theorem}
\emph{(Lebesgue number.) Let $K$ be a compact metric space and
$\{U_\alpha\}$ an open cover of $K$. Then there is a $\delta>0$ --- a
\emph{Lebesgue number} of the cover --- such that every ball
$B(x,\delta)$ with $x \in K$ lies inside a single $U_\alpha$.}
\provenance{Yu Pin, Prop.\ 66 (p.\,112) and the metric version p.\,114.}
\begin{proof}
Suppose no such $\delta$ exists. Then for each $n$ there is $x_n \in K$ with
$B(x_n,1/n)$ contained in no single $U_\alpha$. Compactness gives a
subsequence $x_{n_k} \to x$, and $x$ lies in some $U_\beta$, which contains a
ball $B(x,r)$. For $k$ large enough that $d(x_{n_k},x)<r/2$ and
$1/n_k < r/2$, the triangle inequality gives
$B(x_{n_k},1/n_k) \subset B(x,r) \subset U_\beta$ --- contradicting the
choice of $x_{n_k}$.
\end{proof}

\begin{proofwalk}[How the proof flows]
  \item \textbf{Choose the method.} The claim is that a \emph{single} scale
        works everywhere, and no candidate scale is in sight. So assume none
        exists and let the failure build a sequence --- the standard move
        when a uniform quantity is wanted and only pointwise data is given.
  \item \textbf{Let the counterexamples index themselves.} Failure at scale
        $1/n$ produces a point $x_n$. The sequence is free: it is the
        negation of the conclusion, written out.
  \item \textbf{Spend compactness.} A convergent subsequence appears, and its
        limit lies in some cover element --- which finally supplies a ball of
        a definite radius $r$, the first fixed scale in the argument.
  \item \textbf{Collide the two scales.} For large $k$ the tiny ball around
        $x_{n_k}$ sits inside $B(x,r)$, hence inside one $U_\beta$ ---
        contradicting how $x_{n_k}$ was chosen. Note both smallnesses are
        needed: $d(x_{n_k},x)$ small \emph{and} $1/n_k$ small.
  \item \textbf{Audit.} Compactness is spent once, in step 3; openness of the
        cover once, also in step 3, to get $r$. The lemma is what Rudin's own
        4.19 constructs by hand --- his halving of $\phi(p)$ and taking the
        minimum over a finite subcover produces exactly this $\delta$.
\end{proofwalk}

\begin{unpack}[The lemma hiding inside Rudin's proof of 4.19]
Rudin's uniform continuity proof does not name this, but it manufactures it:
he takes the balls $J(p)$ of radius $\tfrac12\phi(p)$, extracts a finite
subcover, and sets $\delta = \tfrac12\min\phi(p_i)$. That $\delta$ is a
Lebesgue number of the cover $\{J(p)\}$, and the halving trick is exactly the
device that produces one.

Isolate the lemma and 4.19 becomes three lines. Given $\varepsilon>0$, cover
$K$ by the sets $U_x = f^{-1}\bigl(B(f(x),\varepsilon/2)\bigr)$, open by 4.8.
Let $\delta$ be a Lebesgue number. If $d(p,q)<\delta$ then both lie in one
$U_x$, so $f(p)$ and $f(q)$ are both within $\varepsilon/2$ of $f(x)$, hence
within $\varepsilon$ of each other.

The gain is not brevity but reuse: the same lemma proves that compact and
sequentially compact agree, and it is the standard tool whenever a
\emph{single} scale must be extracted from a cover. Yu Pin runs
Heine--Borel itself off it. Keep it as the answer to ``what does compactness
actually hand me?'' --- a uniform scale.
\end{unpack}

% ============================================================
\ssection{Supplement F --- Finer structure: discontinuities and components}

\sitem{S4.20}{Definition}
\emph{(Four kinds of discontinuity.) Let $f$ be discontinuous at $x$. The
discontinuity is}
\begin{enumerate}[label=(\alph*), leftmargin=2.2em, itemsep=2pt, topsep=3pt]
  \item \emph{\textbf{removable} if $f(x+)$ and $f(x-)$ exist and are equal
        (necessarily $\ne f(x)$);}
  \item \emph{\textbf{a jump} if both exist and differ;}
  \item \emph{\textbf{asymptotic} if a one-sided limit fails because $|f|$
        is unbounded near $x$ on that side;}
  \item \emph{\textbf{oscillatory} if a one-sided limit fails while $f$
        stays bounded there.}
\end{enumerate}
\provenance{Tao I, Remark 9.5.4; USTC \cjk{可去间断点/跳跃点}{removable/jump},
bk.\ pp.\,42--43.}

\noindent
Rudin's Definition 4.26 stops at two boxes: (a) and (b) together are his
\emph{simple} discontinuities, and (c) and (d) together are his
\emph{second kind}.\kn{The refinement earns its keep twice. Removable
discontinuities are exactly the ones a change of a single value repairs, so
they are not really properties of the function but of its bookkeeping.
And the asymptotic/oscillatory split is the difference between $1/x$ and
$\sin(1/x)$ at $0$ --- two failures that behave nothing alike, which
Rudin's 4.27(d) illustrates without naming.}

\begin{center}
\begin{tikzpicture}[x=1mm, y=1mm]
  % (a) removable
  \draw[axis] (0,0) -- (34,0);
  \draw[seg] (3,6) -- (16,13);
  \draw[seg] (18,14) -- (31,21);
  \node[ptopen] at (17,13.5) {};
  \node[ptfill, minimum size=2.2pt] at (17,21) {};
  \node[lbl, anchor=north] at (17,-2) {(a) removable};
  % (b) jump
  \begin{scope}[xshift=48mm]
  \draw[axis] (0,0) -- (34,0);
  \draw[seg] (3,7) -- (16,11);
  \node[ptopen] at (17,11.3) {};
  \draw[seg] (18,19) -- (31,23);
  \node[ptfill, minimum size=2.2pt] at (17,19) {};
  \node[lbl, anchor=north] at (17,-2) {(b) jump};
  \end{scope}
  % (c) asymptotic
  \begin{scope}[yshift=-34mm]
  \draw[axis] (0,0) -- (34,0);
  \draw[guide] (17,0) -- (17,23);
  \draw[seg] (5,2.6) .. controls (12,4) and (15.6,9) .. (16.4,22);
  \draw[seg] (17.6,22) .. controls (18.4,9) and (22,4) .. (29,2.6);
  \node[lbl, anchor=north] at (17,-2) {(c) asymptotic};
  \end{scope}
  % (d) oscillatory
  \begin{scope}[xshift=48mm, yshift=-34mm]
  \draw[axis] (0,0) -- (34,0);
  \draw[guide] (17,0) -- (17,22);
  \draw[seg] (3,6.5) .. controls (8,8.5) and (13,10) .. (16.6,11);
  \node[ptopen] at (16.85,11.1) {};
  \draw[seg] plot[smooth] coordinates {
    (17.95,17) (18.05,11) (18.2,5) (18.35,11) (18.55,19)
    (18.8,11) (19.1,4) (19.5,11) (20.0,19.5) (20.6,11)
    (21.3,4) (22.2,11) (23.3,19.5) (24.6,11) (26.1,4.5)
    (27.9,11) (29.6,18.5) (31,14)};
  \node[lbl, anchor=north] at (17,-2) {(d) oscillatory};
  \end{scope}
  % brackets
  \draw[thin] (0,28) -- (0,30) -- (82,30) -- (82,28);
  \node[lbl, anchor=south] at (41,30.6)
    {\emph{first kind} (Rudin: \emph{simple}) --- both one-sided limits exist};
  \draw[thin] (0,-42) -- (0,-44) -- (82,-44) -- (82,-42);
  \node[lbl, anchor=north] at (41,-44.6)
    {\emph{second kind} --- a one-sided limit fails to exist};
\end{tikzpicture}
\end{center}
\figcap{Rudin's two boxes opened into four. His 4.27(c) is case (b) and his
4.27(d) is case (d); cases (a) and (c) he never names. The classification is
per side --- a function may be removable from the left and asymptotic from
the right --- so these are the four pure types, not a partition of all
discontinuities.}

\sitem{S4.21}{Definition}
\emph{$E \subset X$ is \emph{path-connected} if for all $p,q \in E$ there is
a continuous $\gamma : [0,1] \to E$ with $\gamma(0)=p$, $\gamma(1)=q$. Every
path-connected set is connected; the converse is false.}
\provenance{Tao II, Ex.\ 2.4.7 and \S2.4.}

\begin{deeper}[The counterexample is Rudin's own $\sin(1/x)$]
That path-connected implies connected is Rudin 4.22 applied to $\gamma$:
if $E = A \cup B$ were a separation, the sets $\gamma^{-1}(A)$ and
$\gamma^{-1}(B)$ would separate $[0,1]$, which 2.47 forbids.

For the converse, take the closure of
\[ S = \{(x,\sin(1/x)) : 0 < x \le 1\} , \]
which adds the segment $\{0\}\times[-1,1]$. It is connected --- $S$ is
connected as a continuous image of $(0,1]$, and adding limit points of a
connected set preserves connectedness --- but no path joins a point of the
segment to a point of $S$: a path would have to traverse infinitely much
oscillation in finite time, and one shows its first coordinate cannot become
positive. So Rudin's favourite pathological function, introduced at 4.27(d)
purely as a discontinuity, turns out to also separate two notions of
connectedness he never distinguishes.
\end{deeper}

\sitem{S4.22}{Theorem}
\emph{(Components.) Define $p \sim q$ if some connected subset of $X$
contains both. This is an equivalence relation; its classes, the
\emph{connected components} of $X$, are connected and closed, and they
partition $X$.}
\provenance{Tao II, Ex.\ 2.4.6, 2.4.8, 2.4.9.}
\begin{proof}
Reflexivity holds because $\{p\}$ is connected, and symmetry is trivial. For
transitivity, if $A$ is connected containing $p,q$ and $B$ is connected
containing $q,r$, then $A \cup B$ is connected --- a separation of $A\cup B$
would separate one of $A$, $B$, since $q$ lies in both --- and it contains
$p$ and $r$.

The class of $p$ is the union of all connected sets containing $p$; that
union is connected by the same argument, so each class is connected and is
maximal among connected sets meeting it. Finally the closure of a connected
set is connected, so a class equals its own closure and is therefore closed.
\end{proof}

\noindent
Transitivity holds because a union of connected sets with a point in common
is connected, and closedness because the closure of a connected set is
connected.\kn{Rudin never decomposes a space this way, though 2.47 is the
statement that the components of a subset of $\R$ are intervals, and
Exercise 2.29 --- every open set in $\R$ is a countable union of disjoint
segments --- is the component decomposition of an open set. Naming the
notion makes both results one result.}

\sitem{S4.23}{Theorem}
\emph{(Convexity, five ways.) For $f$ on an interval $I$ the following are
equivalent: (i) $f(\lambda x + (1-\lambda)y) \le \lambda f(x) + (1-\lambda)f(y)$
for $\lambda \in (0,1)$; (ii) for $x_1<x_2<x_3$ in $I$,}
\[ \frac{f(x_2)-f(x_1)}{x_2-x_1} \;\le\;
   \frac{f(x_3)-f(x_1)}{x_3-x_1} \;\le\;
   \frac{f(x_3)-f(x_2)}{x_3-x_2}; \]
\emph{(iii) for each fixed $a$, the slope $x \mapsto
\bigl(f(x)-f(a)\bigr)/(x-a)$ is increasing; (iv) the $3\times3$ determinant
with rows $(1,1,1)$, $(x_1,x_2,x_3)$, $\bigl(f(x_1),f(x_2),f(x_3)\bigr)$ is
$\ge 0$; (v) the epigraph $\{(x,y) : y \ge f(x)\}$ is a convex subset of
$\R^2$.}
\provenance{Yu Pin, Thm 106 (pp.\,175--177); ECNU \cjk{凸函数}{convex
functions} \S6.5 gives (i)--(iii).}
\begin{proof}[Proof of (i) $\Leftrightarrow$ (ii)]
Given $x_1<x_2<x_3$, write $x_2 = \lambda x_1 + (1-\lambda)x_3$ with
$\lambda = (x_3-x_2)/(x_3-x_1) \in (0,1)$. Then (i) says
$f(x_2) \le \lambda f(x_1) + (1-\lambda)f(x_3)$, and subtracting $f(x_1)$
from both sides and dividing by $x_2-x_1 = (1-\lambda)(x_3-x_1) > 0$ turns
this into the first inequality of (ii); subtracting instead from $f(x_3)$ and
dividing by $x_3-x_2$ gives the second. Each step is reversible, so (ii)
implies (i).

Statement (iii) is (ii) with $x_1$ held fixed, (iv) is (ii) cleared of
denominators, and (v) says that the chord joining two points of the graph
lies in the region above the graph --- which is (i) verbatim.
\end{proof}

\begin{center}
\begin{tikzpicture}[x=1mm, y=1mm]
  \fill[black!7] (5,19.63) .. controls (16,12.6) and (27,6.6) .. (37,5)
    .. controls (48,5.9) and (61,12.5) .. (70,20.56)
    -- (70,27) -- (5,27) -- cycle;
  \draw[axis] (0,0) -- (72,0);
  \draw[seg] (5,19.63) .. controls (16,12.6) and (27,6.6) .. (37,5)
    .. controls (48,5.9) and (61,12.5) .. (70,20.56);
  \draw[thin] (15,11.91) -- (28,6.16);
  \draw[thin] (15,11.91) -- (58,11.30);
  \draw[thin] (28,6.16) -- (58,11.30);
  \foreach \x/\y in {15/11.91, 28/6.16, 58/11.30}{
    \node[ptfill, minimum size=2.2pt] at (\x,\y) {};
    \draw[guide] (\x,0) -- (\x,\y);}
  \node[mlbl, anchor=north] at (15,-1.4) {$x_1$};
  \node[mlbl, anchor=north] at (28,-1.4) {$x_2$};
  \node[mlbl, anchor=north] at (58,-1.4) {$x_3$};
  \node[lbl, anchor=west] at (34,24) {epigraph: a convex set};
  \node[lbl, anchor=west] at (75,14.5) {chord slopes increase:};
  \node[mlbl, anchor=west] at (75,9.5)
    {$[x_1,x_2] < [x_1,x_3] < [x_2,x_3]$};
\end{tikzpicture}
\end{center}
\figcap{The chord slopes over $[x_1,x_2]$, $[x_1,x_3]$, $[x_2,x_3]$ increase
from left to right --- statement (ii), which is (i) rewritten with
$\lambda = (x_3-x_2)/(x_3-x_1)$. Statement (v) says the same thing without
picking any points at all: everything on or above the graph is a convex set,
which is what ``the chord lies above the graph'' means.}

\begin{deeper}[Why five formulations instead of one]
Rudin's Exercises 4.23--4.24 give the definition and ask for continuity on
the open interval plus the midpoint characterisation. The equivalences above
are what make those exercises routine, and each is the right tool for a
different job. Form (ii) is the one that proves continuity: the slopes from
a fixed interior point are bounded on both sides, so $f$ is Lipschitz on any
compact subinterval of $I^\circ$ --- which incidentally shows why Rudin's
exercise is stated on the \emph{open} interval, since an endpoint can carry
a jump upward and stay convex.

Form (v) is the conceptual one: it converts a statement about a function
into a statement about a \emph{set}, so that convexity of functions and
convexity of sets stop being two ideas. Form (iv) is the computational one,
and form (iii) is what one differentiates in Chapter 5 to get ``$f'$
increasing'' (Theorem 5.11's setting).
\end{deeper}

\begin{pitfall}[Continuity along lines, separately and jointly]
Three conditions on $f : \R^2 \to \R$ at a point, each strictly weaker than
the next:
\begin{enumerate}[label=(\alph*), leftmargin=2.2em, itemsep=1pt, topsep=3pt]
  \item continuous in each variable separately (two lines);
  \item continuous along every straight line through the point;
  \item jointly continuous.
\end{enumerate}
The function $xy/(x^2+y^2)$, with value $0$ at the origin, separates (a)
from (b): it vanishes on both axes but equals $m/(1+m^2)$ along $y=mx$, so
it is separately continuous and not even constant along lines. Rudin's
Exercise 4.7 does the harder separation of (b) from (c): $xy^2/(x^2+y^4)$ is
continuous along every line yet approaches $\tfrac12$ along the parabola
$x=y^2$.

Keeping the hierarchy straight matters in Chapter 9, where partial
derivatives are a statement of type (a) --- which is why 9.21 needs
continuity of the partials, not merely their existence.
\end{pitfall}

\begin{recap}[What the supplements changed]
Take away four things.

\smallskip
\raggedright
\textbf{Criteria produce limits; Rudin's theorems transport them.} S4.1 and
S4.17 are the existence tools, and both need completeness of the target ---
the same property that made Chapter 1 necessary.

\smallskip
\textbf{$o$, $O$, $\sim$ are the working language of limits}, and they are
legal in factors, never in summands (S4.7). Rudin does without them because
he never computes; you will not have that luxury.

\smallskip
\textbf{Compactness is not the only source of good behaviour.} Order does
some of the same work: monotone functions have continuous inverses on any
interval (S4.13), where Rudin's 4.17 needs compactness. Neither theorem
contains the other.

\smallskip
\textbf{Rudin's classifications are coarse on purpose.} Two kinds of
discontinuity become four (S4.20), connectedness sprouts a path-connected
refinement (S4.21), and continuity at a point becomes a \emph{number},
$\omega_f(p)$, which turns ``where is $f$ continuous?'' into a question with
a theorem for an answer (S4.5).
\end{recap}
